% sec6_1.tex — Section 6.1: Modeling Serial Correlation: Linear Processes

Section~2.2 introduced covariance-stationary processes and their autocovariances.
In particular, a (scalar or univariate) \textbf{white noise process} $\{\varepsilon_t\}$ is a zero-mean
covariance-stationary process with no serial correlation:
\[
  \E(\varepsilon_t) = 0, \quad \E(\varepsilon_t^2) = \sigma^2 > 0, \quad
  \E(\varepsilon_t \varepsilon_{t-j}) = 0 \quad \text{for } j \neq 0.
\]
A very important class of covariance-stationary processes, called \textbf{linear processes},
can be created by taking a moving average of a white noise process. This section
studies general properties of linear processes using an apparatus called the \textbf{filter}.

Because the current value of a linear process can depend on possibly infinite
past values of a white noise process, it is convenient to have the white noise process
and the linear process it generates defined for $t = 0, -1, -2, \ldots$ as well
as for $t = 1, 2, \ldots$. Intuitively, a covariance-stationary process defined for all
integers ($t = 0, \pm 1, \pm 2, \ldots$) is a process that started so long ago that its mean and
autocovariances have stabilized to time-invariant constants.

\subsection*{MA(\textit{q})}

The simplest example of linear processes that exhibit serial correlation is a finite-order
moving average process. A process $\{y_t\}$ is called the \textbf{\textit{q}-th order moving
average process} ($\text{MA}(q)$) if it can be written as a weighted average of the current
and most recent $q$ values of a white noise process:
\begin{equation}
  y_t = \mu + \theta_0 \varepsilon_t + \theta_1 \varepsilon_{t-1} + \cdots + \theta_q \varepsilon_{t-q}
  \quad \text{with } \theta_0 = 1.
  \label{eq:6.1.1}
\end{equation}
Evidently, a moving average process is covariance-stationary with mean $\mu$. It is
easy to verify that the $j$-th order autocovariance, $\gamma_j$ ($\equiv \E[(y_t - \mu)(y_{t-j} - \mu)]$), is
\begin{align}
  \gamma_j &= (\theta_j \theta_0 + \theta_{j+1} \theta_1 + \cdots + \theta_q \theta_{q-j}) \sigma^2
         = \sigma^2 \sum_{k=0}^{q-j} \theta_{j+k} \theta_k
         \quad \text{for } j = 0, 1, \ldots, q,
  \label{eq:6.1.2a} \\
  \gamma_j &= 0 \quad \text{for } j > q,
  \label{eq:6.1.2b}
\end{align}
where $\sigma^2 \equiv \E(\varepsilon_t^2)$. (This formula also covers $\gamma_{-j}$ because $\gamma_j = \gamma_{-j}$.)
For example, for $q = 1$, we have
\begin{align*}
  \gamma_0 &= (\theta_0^2 + \theta_1^2) \sigma^2 = (1 + \theta_1^2) \sigma^2 \\
  &\quad \text{(by setting $j = 0$ and $q = 1$ in \eqref{eq:6.1.2a} and noting $\theta_0 = 1$),} \\
  \gamma_1 &= \gamma_{-1} = (\theta_1 \theta_0) \sigma^2 = \theta_1 \sigma^2 \\
  &\quad \text{(by setting $j = 1$ and $q = 1$ in \eqref{eq:6.1.2a} and noting $\theta_0 = 1$),} \\
  \gamma_j &= 0 \quad \text{for } j = \pm 2, \pm 3, \ldots\,.
\end{align*}
The whole profile of autocovariances, $\{\gamma_j\}$, is described by just $q + 1$ parameters
$(\theta_1, \theta_2, \ldots, \theta_q, \text{ and } \sigma^2)$, and the correlogram $\{\rho_j\}$
(where $\rho_j \equiv \gamma_j / \gamma_0$) by $q$ parameters $(\theta_1, \theta_2, \ldots, \theta_q)$.

\subsection*{MA($\infty$) as a Mean Square Limit}

As illustrated in the above example, serial correlation in $\text{MA}(q)$ processes dies out
completely after $q$ lags. Although some time series have this property (we will see
an example in the application below), it is certainly desirable to be able to model
serial correlation that does not have this property. An obvious idea is to have $y_t$
depend on the infinite past by replacing the sum of finitely many terms in \eqref{eq:6.1.1},
\[
  \theta_0 \varepsilon_t + \theta_1 \varepsilon_{t-1} + \cdots + \theta_q \varepsilon_{t-q},
\]
by an infinite sum
\begin{equation}
  \psi_0 \varepsilon_t + \psi_1 \varepsilon_{t-1} + \cdots
  = \sum_{j=0}^{\infty} \psi_j \varepsilon_{t-j}
  \label{eq:6.1.3}
\end{equation}
where $\{\psi_j\}$ is a sequence of real numbers.

For this idea to work, we must make sure that this sum of infinitely many
random variables is well defined, namely, that the partial sum
\begin{equation}
  \sum_{j=0}^{n} \psi_j \varepsilon_{t-j}
  \label{eq:6.1.4}
\end{equation}
converges (in some appropriate sense) to a random variable as $n \to \infty$. One
popular condition under which this happens is that the sequence of real numbers
$\{\psi_j\}$ be \textbf{absolutely summable}:
\begin{equation}
  \sum_{j=0}^{\infty} |\psi_j| < \infty.
  \label{eq:6.1.5}
\end{equation}
(This is a standard short-hand expression in calculus. Equation~\eqref{eq:6.1.5} reads: the
partial sum $\sum_{j=0}^{n} |\psi_j|$ converges to a real number [i.e., a finite limit] as $n \to \infty$
and the infinite sum in \eqref{eq:6.1.5} is \emph{defined} to be this real number.) For the sequence
$\{\psi_j\}$ to be absolutely summable, it is necessary (but certainly not sufficient) that
$\psi_j \to 0$ as $j \to \infty$. Thus, absolute summability requires that the effect of the past
shocks represented by $\psi_j$ eventually die away.

As claimed in the next proposition, the partial sum \eqref{eq:6.1.4} converges in mean
square to some random variable for any given $t$. The mean square limit is unique
in the following sense. If there exists another random variable to which the partial
sum converges in mean square, then (as you will show in Analytical Exercise~1)
the two mean square limits are equal with probability one. We \emph{define} the infinite
sum \eqref{eq:6.1.3} to be the unique mean square limit of the partial sum \eqref{eq:6.1.4} and say
``the infinite sum converges in mean square.''

\begin{proposition}[MA($\infty$) with absolutely summable coefficients]
\label{prop:6.1}
Let $\{\varepsilon_t\}$ be white noise and $\{\psi_j\}$ be a sequence of real numbers that is
absolutely summable. Then
\begin{enumerate}[label=(\alph*)]
\item For each $t$,
\begin{equation}
  y_t = \mu + \sum_{j=0}^{\infty} \psi_j \varepsilon_{t-j}
  \label{eq:6.1.6}
\end{equation}
converges in mean square. $\{y_t\}$ is covariance-stationary. (The process $\{y_t\}$ is
called the \textbf{infinite-order moving-average process} ($\mathrm{MA}(\infty)$).)

\item The mean of $y_t$ is $\mu$. The autocovariances $\{\gamma_j\}$ are given by
\begin{equation}
  \gamma_j = (\psi_j \psi_0 + \psi_{j+1} \psi_1 + \psi_{j+2} \psi_2 + \cdots) \sigma^2
         = \sigma^2 \sum_{k=0}^{\infty} \psi_{j+k} \psi_k
         \quad (j = 0, 1, 2, \ldots).
  \label{eq:6.1.7}
\end{equation}

\item The autocovariances are absolutely summable:
\begin{equation}
  \sum_{j=0}^{\infty} |\gamma_j| < \infty.
  \label{eq:6.1.8}
\end{equation}

\item If, in addition, $\{\varepsilon_t\}$ is i.i.d, then the process $\{y_t\}$ is (strictly) stationary and
ergodic.
\end{enumerate}
\end{proposition}

This result includes $\text{MA}(q)$ processes as a special case with $\psi_j = \theta_j$ for
$j = 0, 1, \ldots, q$ and $\psi_j = 0$ for $j > q$. Evidently, this sequence $\{\psi_j\}$ is absolutely
summable. Proving parts (a)--(c) is Analytical Exercise~2. Proving (a) would use
the fact from calculus that Cauchy sequences converge (see Review Question~1
about Cauchy sequences). Part~(b) is an obvious extrapolation to infinity of the
corresponding results on the mean and autocovariances of finite-order MA processes,
but the extrapolation involves interchanging the order of expectations and
summations. For example, to prove that the mean of the $\text{MA}(\infty)$ process is $\mu$, you
would have to show that the expected value of the mean square limit of \eqref{eq:6.1.3} is
the limit of the expected value of \eqref{eq:6.1.4}. The fact that this operation is legitimate
under mean square convergence would be used in the proof of part~(b). Part~(c)
says, intuitively, that if the effect of past shocks represented by $\psi_j$ dies out fast
enough to make $\{\psi_j\}$ absolutely summable, then serial correlation dies out just as
fast. Absolute summability of autocovariances will play a key role in the theory of
covariance-stationary processes to be developed in this chapter. For a proof of~(d),
see, e.g., Hannan (1970, p.~204, Theorem~3).\footnote{Absolute summability of $\{\psi_j\}$ is more than enough for part~(a) to hold. As shown in Analytical Exercise~2, square summability, which is weaker than absolute summability, is enough to ensure that the infinite sum is well defined as a mean square limit. However, for the other parts of this proposition to hold, absolute summability is needed.}

The $\text{MA}(\infty)$ process in \eqref{eq:6.1.6} is said to be \textbf{one-sided} because the moving
average does not include future values of $\varepsilon_t$. It is a special case of a \textbf{linear process},
which is defined to be two-sided moving averages of the form:
\[
  y_t = \mu + \sum_{j=-\infty}^{\infty} \psi_j \varepsilon_{t-j}
  \quad \text{with} \quad \sum_{j=-\infty}^{\infty} |\psi_j| < \infty.
\]
The proposition (and all the results in this section) can easily be extended to general
linear processes. We nevertheless focus on one-sided moving averages, only
because two-sided moving averages are very rarely encountered in economics.

Proposition~6.1 can be generalized to the case where the process $\{y_t\}$ is a moving
average of a general covariance-stationary process rather than of a white noise
process. In particular, absolute summability of autocovariances survives the operation
of taking an absolutely summable weighted average.

\begin{proposition}[Filtering covariance-stationary processes]
\label{prop:6.2}
Let $\{x_t\}$ be a covariance-stationary process and $\{h_j\}$ be a sequence of real numbers
that is absolutely summable. Then
\begin{enumerate}[label=(\alph*)]
\item For each $t$, the infinite sum
\begin{equation}
  y_t = \sum_{j=0}^{\infty} h_j x_{t-j}
  \label{eq:6.1.9}
\end{equation}
converges in mean square. The process $\{y_t\}$ is covariance-stationary.

\item If, furthermore, the autocovariances of $\{x_t\}$ are absolutely summable, then so
are the autocovariances of $\{y_t\}$.
\end{enumerate}
\end{proposition}

We defer to, e.g., Fuller (1996, Theorem~4.3.1) for a proof, but the basic technique
is the same as in the proof of Proposition~6.1. Deriving the autocovariances of $\{y_t\}$,
which would generalize \eqref{eq:6.1.7}, is Analytical Exercise~3.

\subsection*{Filters}

The operation of taking a weighted average of (possibly infinitely many) successive
values of a process, as in \eqref{eq:6.1.6} and \eqref{eq:6.1.9}, is called \textbf{filtering}. It can be expressed
compactly if we use the device called the \textbf{lag operator} $L$, defined by the relation
$L^j x_t = x_{t-j}$. We now introduce some concepts associated with the lag operator
that will be useful in time-series analysis.\footnote{For a fuller treatment of filters, see, e.g., Gourieroux and Monfort (1997, Sections~5.2 and~7.4).}

\subsection*{Filters Defined}

For a given arbitrary sequence of real numbers, $(\alpha_0, \alpha_1, \alpha_2, \ldots)$, define a
\textbf{filter} by
\begin{equation}
  \alpha(L) \equiv \alpha_0 + \alpha_1 L + \alpha_2 L^2 + \cdots\,.
  \label{eq:6.1.10}
\end{equation}
If we just mechanically and mindlessly apply the definition of the lag operator to
an input process $\{x_t\}$, we get
\begin{equation}
  \alpha(L) x_t = \alpha_0 x_t + \alpha_1 L x_t + \alpha_2 L^2 x_t + \cdots
               = \alpha_0 x_t + \alpha_1 x_{t-1} + \alpha_2 x_{t-2} + \cdots
               = \sum_{j=0}^{\infty} \alpha_j x_{t-j}.
  \label{eq:6.1.11}
\end{equation}
This is indeed the definition of $\alpha(L) x_t$. If $x_t = c$ (a constant), $\alpha(L) c = \alpha(1) c =
c \sum_{j=0}^{\infty} \alpha_j$. If $\alpha_j \neq 0$ for $j = p$ and $\alpha_j = 0$ for $j > p$, the filter reduces to a \textbf{\textit{p}-th
degree lag polynomial}:
\[
  \alpha(L) = \alpha_0 + \alpha_1 L + \cdots + \alpha_p L^p.
\]
When applied to an input process, it creates a weighted average of the current
and $p$ most recent values of the process. Proposition~6.2 assures us that the object
$\alpha(L) x_t$ defined by \eqref{eq:6.1.11} is a well-defined random variable forming a covariance-stationary process if the sequence $\{\alpha_j\}$ is absolutely summable and if the input process
$\{x_t\}$ is covariance-stationary. A filter with absolutely summable coefficients
will be referred to (in this book) as an \textbf{absolutely summable filter}. To use some
fancy set theory language, an absolutely summable filter is a mapping from the set
of covariance-stationary processes to itself.

\subsection*{Products of Filters}

Let $\{\alpha_j\}$ and $\{\beta_j\}$ be two arbitrary sequences of real numbers and define the sequence
$\{\delta_j\}$ by the relation
\begin{equation}
\begin{aligned}
  \alpha_0 \beta_0 &= \delta_0, \\
  \alpha_0 \beta_1 + \alpha_1 \beta_0 &= \delta_1, \\
  \alpha_0 \beta_2 + \alpha_1 \beta_1 + \alpha_2 \beta_0 &= \delta_2, \\
  &\;\;\vdots \\
  \alpha_0 \beta_j + \alpha_1 \beta_{j-1} + \alpha_2 \beta_{j-2} + \cdots + \alpha_{j-1} \beta_1 + \alpha_j \beta_0 &= \delta_j, \\
  &\;\;\vdots\,.
\end{aligned}
\label{eq:6.1.12}
\end{equation}
The sequence $\{\delta_j\}$ created from this convoluted formula is called the \textbf{convolution}
of $\{\alpha_j\}$ and $\{\beta_j\}$. Let $\alpha(L) \equiv \alpha_0 + \alpha_1 L + \alpha_2 L^2 + \cdots$,
$\beta(L) \equiv \beta_0 + \beta_1 L + \beta_2 L^2 + \cdots$,
and $\delta(L) = \delta_0 + \delta_1 L + \delta_2 L^2 + \cdots$ be the associated filters. We define the
\textbf{product} of two filters $\alpha(L)$ and $\delta(L)$ and write it as
\[
  \delta(L) = \alpha(L)\,\beta(L).
\]
For example, for $\alpha(L) = 1 + \alpha_1 L$ and $\beta(L) = 1 + \beta_1 L$, the convolution formula
yields
\[
  (1 + \alpha_1 L)(1 + \beta_1 L) = 1 + (\alpha_1 + \beta_1) L + \alpha_1 \beta_1 L^2.
\]
(If you are familiar with the product of two polynomials of the usual sort from calculus,
you can immediately see that this definition for filters is strictly analogous.)
As is clear from the definition, filters are \textbf{commutative}:
\[
  \alpha(L)\,\beta(L) = \beta(L)\,\alpha(L).
\]
Commutativity, however, will not carry over to matrix filters (see Section~6.3
below).

If the two sequences $\{\alpha_j\}$ and $\{\beta_j\}$ are absolutely summable and if $\{x_t\}$ is
covariance-stationary, then Proposition~6.2 assures us that $\{\beta(L) x_t\}$ is covariance-stationary and so $\alpha(L)\,\beta(L) x_t$ is a well-defined random variable, equal to $\delta(L) x_t$.
Thus, if $\alpha(L)$ and $\beta(L)$ are absolutely summable and $\{x_t\}$ is covariance-stationary,
then
\begin{equation}
  \text{``}\alpha(L)\,\beta(L) = \delta(L)\text{''} \;\Rightarrow\;
  \text{``}\alpha(L)\,\beta(L) x_t = \delta(L) x_t.\text{''}
  \label{eq:6.1.13}
\end{equation}
(Incidentally, $\{\delta_j\}$ is absolutely summable [see, e.g., Fuller, 1996, p.~30].)

\subsection*{Inverses}

Of particular interest is the case when $\delta(L) = 1$ (a very special filter) so that the
two filters $\alpha(L)$ and $\beta(L)$ satisfy $\alpha(L)\,\beta(L) = 1$. We say that $\beta(L)$ is the \textbf{inverse}
of $\alpha(L)$ and denote it as $\alpha(L)^{-1}$ or $1/\alpha(L)$. That is,
\begin{equation}
  \alpha(L)^{-1} \text{ is a filter satisfying } \alpha(L)\,\alpha(L)^{-1} = 1.
  \label{eq:6.1.14}
\end{equation}
As long as $\alpha_0 \neq 0$ in $\alpha(L) = \alpha_0 + \alpha_1 L + \cdots$, the inverse of $\alpha(L)$ can be defined
for any arbitrary sequence $\{\alpha_j\}$ because the set of equations \eqref{eq:6.1.12} with $\delta_0 = 1$,
$\delta_j = 0$ ($j = 1, 2, \ldots$) can be solved successively for $\{\beta_j\}$:
\[
  \beta_0 = \frac{1}{\alpha_0}, \quad
  \beta_1 = -\frac{\alpha_1 \beta_0}{\alpha_0} = -\frac{\alpha_1}{\alpha_0^2}, \quad \text{etc.}
\]
For example, consider the filter $1 - L$. Its inverse is given by
\[
  (1 - L)^{-1} = 1 + L + L^2 + L^3 + \cdots\,.
\]
The coefficient sequence is obviously not absolutely summable. This example
illustrates the point that an inverse filter may not be absolutely summable.

It is left as Review Question~3 to show that
\begin{align}
  \alpha(L)\,\alpha(L)^{-1} &= \alpha(L)^{-1}\,\alpha(L)
  \quad \text{(so inverses are commutative),}
  \label{eq:6.1.15a} \\[4pt]
  \text{``}\alpha(L)\,\beta(L) = \delta(L)\text{''} &\;\Leftrightarrow\;
  \text{``}\beta(L) = \alpha(L)^{-1}\,\delta(L)\text{''} \;\Leftrightarrow\;
  \text{``}\alpha(L) = \delta(L)\,\beta(L)^{-1},\text{''}
  \label{eq:6.1.15b}
\end{align}
provided $\alpha_0 \neq 0$ and $\beta_0 \neq 0$. Therefore, for example, to solve $\alpha(L)\,\beta(L) = \delta(L)$
for $\beta(L)$, you ``multiply by the left'' by the inverse of $\alpha(L)$. We have noted above
that commutativity does not necessarily hold for matrix filters. As will be shown
in Section~6.3, for inverses, commutativity does generalize to matrix filters.

\subsection*{Inverting Lag Polynomials}

In the next section on ARMA processes, it will become necessary to calculate the
inverse of a $p$-th degree lag polynomial $\phi(L)$
\[
  \phi(L) = 1 - \phi_1 L - \phi_2 L^2 - \cdots - \phi_p L^p.
\]
Since $\phi_0 = 1 \neq 0$, the inverse can be defined. Let $\psi(L) \equiv \phi(L)^{-1}$. By definition,
it satisfies
\[
  \phi(L)\,\psi(L) = 1.
\]
The convolution formula \eqref{eq:6.1.12} provides an algorithm for solving this for $\psi(L)$.
Setting $\alpha(L) = \phi(L)$, $\beta(L) = \psi(L)$, and $\delta(L) = 1$ in the formula yields
\begin{equation}
\begin{alignedat}{5}
  \text{constant:}&\quad & \psi_0 && && &= 1, \\
  L\text{:}&\quad & \psi_1 &\;-\; \phi_1 \psi_0 && &= 0, \\
  L^2\text{:}&\quad & \psi_2 &\;-\; \phi_1 \psi_1 &\;-\; \phi_2 \psi_0 &= 0, \\
  &\quad & &\;\;\vdots && \\
  L^{p-1}\text{:}&\quad & \psi_{p-1} &\;-\; \phi_1 \psi_{p-2} - \phi_2 \psi_{p-3} - \cdots - \phi_{p-1} \psi_0 &\quad &= 0, \\
  L^p\text{:}&\quad & \psi_p &\;-\; \phi_1 \psi_{p-1} - \phi_2 \psi_{p-2} - \cdots - \phi_{p-1} \psi_1 - \phi_p \psi_0 &\quad &= 0, \\
  L^{p+1}\text{:}&\quad & \psi_{p+1} &\;-\; \phi_1 \psi_p \;-\; \phi_2 \psi_{p-1} - \cdots - \phi_{p-1} \psi_2 - \phi_p \psi_1 &\quad &= 0, \\
  L^{p+2}\text{:}&\quad & \psi_{p+2} &\;-\; \phi_1 \psi_{p+1} - \phi_2 \psi_p \;- \cdots - \phi_{p-1} \psi_3 - \phi_p \psi_2 &\quad &= 0, \\
  &\quad & &\;\;\vdots && &\;\;.
\end{alignedat}
\label{eq:6.1.16}
\end{equation}
These equations can easily be solved successively for $(\psi_0, \psi_1, \psi_2, \ldots)$ as
\[
  \psi_0 = 1, \quad \psi_1 = \phi_1, \quad \psi_2 = \phi_2 + \phi_1^2, \quad \text{etc.}
\]
Also, notice that for sufficiently large $j$ (actually for $j \geq p$), $\{\psi_j\}$ follows the $p$-th
order homogeneous difference equation
\begin{equation}
  \psi_j - \phi_1 \psi_{j-1} - \phi_2 \psi_{j-2} - \cdots - \phi_{p-1} \psi_{j-p+1} - \phi_p \psi_{j-p} = 0.
  \label{eq:6.1.17}
\end{equation}
So, once the first $p$ coefficients $(\psi_0, \psi_1, \ldots, \psi_{p-1})$ are calculated, we can use this
homogeneous $p$-th order difference equation to generate the rest of the coefficients
$(\psi_j, j \geq p)$ with those first $p$ coefficients as the initial condition.

As we know from the theory of homogeneous difference equations (summarized
in Analytical Exercise~4), the solution sequence $\{\psi_j\}$ to \eqref{eq:6.1.17} eventually
starts declining at a geometric rate if what is known as the \textbf{stability condition}
holds. The condition states:

All the roots of the $p$-th degree polynomial equation in $z$
\begin{equation}
  \phi(z) = 0 \quad \text{where } \phi(z) \equiv 1 - \phi_1 z - \phi_2 z^2 - \cdots - \phi_p z^p
  \label{eq:6.1.18}
\end{equation}
are greater than 1 in absolute value.

The roots of the polynomial equation $\phi(z) = 0$ can be complex numbers. If your
command of complex numbers is rusty, a complex number can be written as
\[
  z = a + bi,
\]
where $a$ and $b$ are two real numbers and $i = \sqrt{-1}$. So a complex number can
be represented as a point in a two-dimensional diagram with the horizontal axis
measuring $a$ (the real part) and the vertical axis measuring $b$ (the imaginary part).
Its absolute value is $|z| = \sqrt{a^2 + b^2}$, which is the distance between the origin and
the point representing $z$. We say that $z$ lies outside the unit circle when $|z| > 1$.
So the stability condition above can be restated as requiring that the roots lie
outside the unit circle. Because $z = \lambda$ solves $\phi(z) = 0$ whenever $z = 1/\lambda$ solves
$z^p - \phi_1 z^{p-1} - \phi_2 z^{p-2} - \cdots - \phi_{p-1} z - \phi_p = 0$, the stability condition can be
equivalently stated as

All the roots of the $p$-th order polynomial equation
\begin{equation}
  z^p - \phi_1 z^{p-1} - \phi_2 z^{p-2} - \cdots - \phi_{p-1} z - \phi_p = 0
  \label{eq:6.1.18prime}
  \tag{6.1.18$'$}
\end{equation}
are less than 1 in absolute value (i.e., lie inside the unit circle).

Under this condition, as you will be asked to prove (Analytical Exercise~4),
there exist two real numbers, $A > 0$ and $0 < b < 1$, such that
\begin{equation}
  |\psi_j| < A b^j \quad \text{for all } j.
  \label{eq:6.1.19}
\end{equation}
Given \eqref{eq:6.1.19}, the absolute summability of $\{\psi_j\}$ follows because
\[
  \sum_{j=0}^{\infty} |\psi_j| < \sum_{j=0}^{\infty} A b^j = \frac{A}{1 - b} < \infty.
\]
Thus, we have proved

\begin{proposition}[Absolutely summable inverses of lag polynomials]
\label{prop:6.3}
Consider a $p$-th degree lag polynomial $\phi(L) = 1 - \phi_1 L - \phi_2 L^2 - \cdots - \phi_p L^p$,
and let $\psi(L) \equiv \phi(L)^{-1}$. If the associated $p$-th degree polynomial equation
$\phi(z) = 0$ satisfies the stability condition \eqref{eq:6.1.18} or \eqref{eq:6.1.18prime}, then the coefficient
sequence $\{\psi_j\}$ of $\psi(L)$ is bounded in absolute value by a geometrically declining
sequence (as in \eqref{eq:6.1.19}) and hence is absolutely summable.
\end{proposition}

To illustrate, consider a lag polynomial of degree one, $1 - \phi L$. The root of the
associated polynomial equation $1 - \phi z = 0$ is $1/\phi$. The stability condition is that
$|1/\phi| > 1$ or $|\phi| < 1$. The filter
\begin{equation}
  1 + \phi L + \phi^2 L^2 + \cdots
  \label{eq:6.1.20}
\end{equation}
is evidently its inverse. The coefficient sequence $\{\phi^j\}$ is indeed bounded in absolute
value by a geometrically declining sequence (for example, set $A = 1.1$ and $b = \phi$).
